% AUTHOR'S CLASSROOM HANDOUT: 1.3 Bosons (six pages), prepared 2024-01-11.
% Rewritten as a self-contained textbook derivation; the Chapter 6 normalization is reused.
\section{Bosons: Fock Space and Creation and Annihilation Operators}

A completely symmetric $N$-particle state can be labeled by the \emph{occupation numbers} $n_i$ of one-particle orbitals $\ket i$. The actual particle labels have disappeared from this description. Define
\begin{equation}
\ket{n_1,n_2,\ldots}_{\mathrm B},\qquad
n_i=0,1,2,\ldots,\qquad\sum_i n_i=N,
\end{equation}
by the normalized symmetrization in the preceding chapter. Because different occupation lists are orthogonal,
\begin{equation}
{}_{\mathrm B}\braket{n_1,n_2,\ldots}{n_1',n_2',\ldots}_{\mathrm B}
=\prod_i\delta_{n_i n_i'}.
\end{equation}
For a fixed $N$ these states span the symmetric $N$-particle Hilbert space $\mathcal H_N^{(\mathrm B)}$. Combining the spaces for every particle number leads to the Bose Fock space,
\begin{equation}
\boxed{\mathcal F_{\mathrm B}=\bigoplus_{N=0}^\infty\mathcal H_N^{(\mathrm B)}.}
\end{equation}
The $N=0$ subspace contains a distinguished normalized state, the \emph{vacuum}, written $\ket0=\ket{0,0,\ldots}$. A state with several occupied orbitals belongs to one fixed-$N$ sector; a general Fock-space state may be a superposition of sectors.

\subsection*{1) Creation and annihilation operators}

The occupation-number basis makes the effect of adding or removing one particle in orbital $i$ explicit. Define
\begin{equation}
\boxed{\begin{aligned}
\op a_i^\dagger\ket{\ldots,n_i,\ldots}_{\mathrm B}
  &=\sqrt{n_i+1}\ket{\ldots,n_i+1,\ldots}_{\mathrm B},\\
\op a_i\ket{\ldots,n_i,\ldots}_{\mathrm B}
  &=\sqrt{n_i}\ket{\ldots,n_i-1,\ldots}_{\mathrm B}.
\end{aligned}}
\end{equation}
The second formula gives zero when $n_i=0$. The square-root factors follow from the normalized symmetric states, including the factorial $n_i!$. In particular,
\begin{equation}
\op a_i\ket0=0,\qquad\ket i=\op a_i^\dagger\ket0,\qquad
\ket{n_1,n_2,\ldots}_{\mathrm B}
=\prod_i\frac{(\op a_i^\dagger)^{n_i}}{\sqrt{n_i!}}\ket0.
\end{equation}
The order of factors is immaterial for bosons. To check the normalization of the last expression, use the norm relation
\begin{equation}
\bigl\|\op a_i^\dagger\ket{n_i}\bigr\|^2=n_i+1
\quad\Longrightarrow\quad
\bigl\|(\op a_i^\dagger)^{n_i}\ket0\bigr\|^2=n_i!.
\end{equation}
Thus the operator algebra reproduces the explicit permutation sum without referring to labels of individual particles.

\subsection*{2) Derivation of Bose commutation relations}

For different modes $i\ne j$, creation of particles in modes $i,j$ can be carried out in either order. Applying the preceding formulas twice yields
\begin{equation}
\op a_i\op a_j\ket{\{n_k\}}
  =\sqrt{n_i n_j}\ket{\ldots,n_i-1,\ldots,n_j-1,\ldots}
  =\op a_j\op a_i\ket{\{n_k\}}.
\end{equation}
The same argument applies to two creation operators and to $\op a_i\op a_j^\dagger$ for $i\ne j$. For one fixed mode,
\begin{equation}
\begin{aligned}
\op a_i\op a_i^\dagger\ket{\{n_k\}}
 &=(n_i+1)\ket{\{n_k\}},\\
\op a_i^\dagger\op a_i\ket{\{n_k\}}
 &=n_i\ket{\{n_k\}}.
\end{aligned}
\end{equation}
Their difference is the identity on every occupation-number basis state. We have therefore established
\begin{equation}
\boxed{[\op a_i,\op a_j]=0,\quad
[\op a_i^\dagger,\op a_j^\dagger]=0,\quad
[\op a_i,\op a_j^\dagger]=\delta_{ij}\mathbf1.}
\end{equation}
Here $[\op A,\op B]=\op A\op B-\op B\op A$. The positive sign in the last relation is characteristic of bosonic statistics.

\subsection*{3) Particle-number operators and independent particles}

Define the occupation-number operator for mode $i$ by
\begin{equation}
\op n_i=\op a_i^\dagger\op a_i,\qquad
\op n_i\ket{\{n_k\}}=n_i\ket{\{n_k\}}.
\end{equation}
Summing over all modes gives the total particle-number operator,
\begin{equation}
\boxed{\op N=\sum_i\op n_i=\sum_i\op a_i^\dagger\op a_i.}
\end{equation}
For a noninteracting system with one-particle energies $\epsilon_i$, the many-particle Hamiltonian acts on an occupation state as
\begin{equation}
\op H_0\ket{\{n_i\}}=\left(\sum_i n_i\epsilon_i\right)\ket{\{n_i\}},
\end{equation}
which is reproduced exactly by
\begin{equation}
\boxed{\op H_0=\sum_i\epsilon_i\op a_i^\dagger\op a_i.}
\end{equation}
This formula is a first example of a sum over individually identical one-particle operators represented entirely in terms of occupation modes.

\subsection*{4) General one-particle operator}

Consider a one-particle observable $\op t$ and the additive many-particle operator
\begin{equation}
\op{\mathcal O}=\sum_{\alpha=1}^N\op t_\alpha,\qquad
 t_{ij}=\mel i{\op t}j.
\end{equation}
The completeness of the one-particle basis gives
\begin{equation}
\op t=\sum_{ij}t_{ij}\ket i\bra j,\qquad
\op{\mathcal O}=\sum_{ij}t_{ij}\left(\sum_{\alpha=1}^N\ket i_\alpha\bra j_\alpha\right).
\end{equation}
The operator in parentheses removes one particle from $\ket j$ and returns one in $\ket i$. On a bosonic occupation state with $i\ne j$, there are $n_j$ possible removals and a corresponding symmetrization factor, which together produce
\begin{equation}
\begin{aligned}
\left(\sum_{\alpha}\ket i_\alpha\bra j_\alpha\right)
 \ket{\ldots,n_i,\ldots,n_j,\ldots}_{\mathrm B}
 &=\sqrt{n_j(n_i+1)}
 \ket{\ldots,n_i+1,\ldots,n_j-1,\ldots}_{\mathrm B}\\
 &=\op a_i^\dagger\op a_j\ket{\{n_k\}}_{\mathrm B}.
\end{aligned}
\end{equation}
For $i=j$ it simply counts $n_i$ particles. Since the actions agree on a complete Fock basis,
\begin{equation}
\boxed{\op{\mathcal O}=\sum_{ij}t_{ij}\op a_i^\dagger\op a_j.}
\end{equation}
The noninteracting Hamiltonian above follows by setting $t_{ij}=\epsilon_i\delta_{ij}$. Notice the operator order: first the particle is removed from $j$, and then one is created in $i$.

\subsection*{5) Two-particle interactions}

For a symmetric pair potential $V(x_1,x_2)=V(x_2,x_1)$, the first-quantized operator counts each unordered pair once:
\begin{equation}
\op V=\sum_{\alpha<\beta}V(\op x_\alpha,\op x_\beta)
 =\frac12\sum_{\alpha\ne\beta}V(\op x_\alpha,\op x_\beta).
\end{equation}
Introduce the \emph{unsymmetrized} orbital matrix elements
\begin{equation}
V_{ij;kl}=\int\dd x\,\dd y\,
  \varphi_i^*(x)\varphi_j^*(y)V(x,y)\varphi_k(x)\varphi_l(y).
\end{equation}
They obey $V_{ij;kl}=V_{ji;lk}$ for a symmetric interaction. Acting on an occupation state with the ordered product $\op a_l\op a_k$ removes the two incoming particles, and $\op a_i^\dagger\op a_j^\dagger$ inserts the two outgoing particles. Summing over the orbital indices exactly reconstructs the matrix elements of $\sum_{\alpha\ne\beta}V_{\alpha\beta}$; the factor $1/2$ removes double counting. Thus
\begin{equation}
\boxed{\op V=\frac12\sum_{ijkl}V_{ij;kl}
\op a_i^\dagger\op a_j^\dagger\op a_l\op a_k.}
\end{equation}
The ordering of $k,l$ on the right is deliberate and will be retained for fermions. This formula uses unsymmetrized matrix elements; antisymmetrizing those matrix elements would change the appropriate prefactor and indexing conventions.

\interactiveguide{IC-06-03}{Bosonic Fock occupations}{After constructing the bosonic ladder operators, apply creation and annihilation to different number states. Check the $sqrt{n+1}$ and $sqrt n$ factors by evaluating state norms and commutators.}
